\honest{The Affect Heuristic}{Eliezer Yudkowsky}{2007}

The affect heuristic is judging by a quick feeling of good or bad. It produces some of my favorite biases, and I take them in rising order of craziness.

First, the clock. Asked what they would pay to insure a clock for \$100 in shipping, people offered more than twice as much when the clock was a gift from their grandparents than when it came from a remote relative. The payout is the same for either clock. You could say they were buying consolation, but I move on. \nb{Consolation is the explanation Hsee and Kunreuther themselves proposed: their ``consolation hypothesis'' holds that people see insurance compensation as a token of consolation.}

Second, Yamagishi's disease that kills 1,286 people out of every 10,000, judged more dangerous than one that is ``24.14\% likely to be fatal.'' ``Apparently,'' a thousand dead bodies are more alarming than ``a single person who's more likely to survive than not.'' \nb{In summaries of the study, both figures were death rates in a population, such as 24.14 out of 100 people, and later work classes the result with the ratio bias. The ``single person'' is the post's own framing.}

Third, an airport measure that would save 150 lives drew a mean support of 10.4 out of 20, and one that would save 98\% of 150 lives drew 13.6. Fourth, people drawing jelly beans for \$1 often preferred 7 red beans in 100 to 1 in 10, and said that although they knew the odds were against them, they felt luckier with more red beans. With irony: this ``makes perfect sense.''

As I discussed in ``The Scales of Justice, the Notebook of Rationality,'' Finucane and colleagues found that information about benefits lowered perceived risks, information about risks lowered perceived benefits, ``and so on across the quadrants.'' Time pressure ``greatly'' increased the effect. \nb{``Greatly'' is also the authors' word in later summaries. In the 2000 paper, a study of 54 students, the mean correlation between risk and benefit ratings was $-0.45$ under time pressure and $-0.33$ without it, a difference that ``approached significance.''} Ganzach found that analysts judged unfamiliar stocks as all good or all bad, while for familiar stocks risk and return went together, as economic theory predicts.

For more, I recommend Slovic's ``Rational Actors or Rational Fools.'' \nb{That article covers the airport and Finucane studies. The others are summarized in the same authors' ``The Affect Heuristic.''}
